\section{Grundlagen}

\subsection{Gleichungen zur Modellbildung}
	Kräftegleichungen:
	\begin{mdframed}[style=exercise]
		Federkraft: \(F_F\) = \(k_F \cdot x\)\\
		Dampfkraft: \(F_D\) = \(k_D \cdot v\) = \(k_D \cdot \dot{x}\)\\
		Trägheitskraft: \(F_{Tr}\) = \(m\cdot a\) = \(m\cdot \ddot{x}\)\\
		Erdanziehungskraft: \(F_G\) = \(m\cdot g\)
	\end{mdframed}
	Moment Gleichungen:
	\begin{mdframed}[style=exercise]
		Widerstandsmoment: \(M_w = k_D \cdot \omega\)\\
		Trägheitsmoment: \(M_{TR} = J \cdot \dot{\omega}\)
	\end{mdframed}
	Spannungsgleichung:
	\begin{mdframed}[style=exercise]
		\[
			U = L\cdot \frac{di}{dt}+i\cdot R+\frac{1}{C} \cdot \int i
		\]
	\end{mdframed}

	Rotation in Flüssigkeit:
	\[
		M=M_{\texttt{Träg}}+M_{\texttt{Brems}}=J \cdot \dot{\omega} +k_{\texttt{Flüssigkeit}} \cdot \omega
	\]

\subsection{Mathematische Grundlagen}
	Für kleine Winkel \(\alpha\) gilt: sin(\(\alpha\)) = \(\alpha\)\\
	\\
	Partialbruchzerlegung (siehe Papula s.157ff)\\
	\\
	Anfangswertsatz: \(x(t \rightarrow +0)=\lim _{s \rightarrow \infty} s \cdot X(s)\)\\
	Endwertsatz: \(x(t \rightarrow \infty)=\lim _{s \rightarrow 0} s \cdot X(s)\)

\subsection{Linearisierung um Arbeitspunkt}
	\begin{align*}
		x_{a}(t) &= x_{a,AP} + \Delta x_{a}(t) \\ 
		&\approx f(x_{e1,AP}, ..., x_{en,AP}) + \sum_{m=1}^{n} \left. \frac{\partial f}{\partial x_{em}} \right|_{AP} \cdot \Delta x_{em}(t)
	\end{align*}

\subsection{Eigenschaften von LTI-Systemen}
	\begin{mdframed}[style=exercise]
		\begin{enumerate}
			\item BIBO-Stabilität {\footnotesize (Bounded Input / Bounded Output)}\\
				\(\abs{y(t)} < \infty \qquad \text{für} \quad \abs{x(t)} < \infty\)
			\item Linearität\\
				\(f\qty{\sum_{k=1}^{N}a_n x_n(t)}=\sum_{n=1}^{N}f\qty{a_n x_n(t)}\)
			\item Zeitinvarianz\\
				\(f\qty{x(t-t_0)} =y(t-t_0)\)
			\item Kausalität\\
				\(x(t)=0 \; , \; y(t)=0 \qquad \text{für} \quad t < 0\)
		\end{enumerate}
	\end{mdframed}

\subsection{Systemantwort}
	\begin{mdframed}[style=exercise]
		\textbf{Sprungantwort} \((x = \sigma)\)
		\begin{align*}
			g(t) &= \mathcal{L}^{-1}\qty{G(s)} = \mathcal{L}^{-1}\qty{\frac{F(s)}{s}}\\
			g(t) &= \int_{0}^{t} h(\tau) \dd{\tau}
		\end{align*}
	\end{mdframed}
	\begin{mdframed}[style=exercise]
		\textbf{Impulsantwort} \((x = \delta)\)
		\begin{align*}
			h(t) &= \int_{-\infty}^{\infty} x(\tau) \delta(t-\tau) \dd{\tau}\\
			h(t) &= \mathcal{L}^{-1}\qty{F(s)}\\
			h(t) &= \frac{d}{\dd{t}} g(t)
		\end{align*}
	\end{mdframed}

\subsection{Abtasttheorem}
	Durch die Abtastung wird das Spektrum von \(f(t)\) unendlich oft um die Frequenzen \(n\cdot \omega_a\) reproduziert.
	\begin{mdframed}[style=exercise]
		\begin{align}
			F_A(\omega) &= \frac{1}{T_A} \sum_{n=-\infty}^{\infty} F(\omega-n\omega_A)\\
			2\omega_g &\leq \omega_A
		\end{align}
	\end{mdframed}